\chapter*{質疑応答}

\note{卒研発表後に，発表当日の質疑で提起された質問とそれに対する答えをまとめてください．}

\begin{table}[h]
\centering
\large
\begin{tabular}{ll}

& 伊藤　雄一　情報テクノロジー学科　教授 \\ \hline
\multicolumn{1}{|l|}{Q} & \multicolumn{1}{l|}{\begin{tabular}[c]{@{}l@{}}\parbox{37zw}{\strut{}
今後より高度なGPTが発表された場合でも，提案法は有効と考えられるか？あるいは何か改善すべき点があるか？
\strut}\end{tabular}} \\ \hline
\multicolumn{1}{|l|}{A} & \multicolumn{1}{l|}{\begin{tabular}[c]{@{}l@{}}\parbox{37zw}{\strut{}
提案法では，XXXという考え方を利用しているため，どんなにGPTが高度になっても対応できると考えられる．
\strut}\end{tabular}} \\ \hline

\end{tabular}
\end{table}

\begin{table}[h]
\centering
\large
\begin{tabular}{ll}

& Martin J. D\"{u}rst 情報テクノロジー学科　教授 \\ \hline
\multicolumn{1}{|l|}{Q} & \multicolumn{1}{l|}{\begin{tabular}[c]{@{}l@{}}\parbox{37zw}{\strut{}
「ChatGPT と人を見分ける」ということがタイトルにあるが，評価実験では ChatGPT によって生成された文章だけを入力対象にしている．人が書いた文章を入力にしてそれが人によって書かれたものであると判定できることを確認しなくて良いか？極端な話，毎回ランダムに 95\% の確率で ChatGPT と答えれば，評価実験と同じ結果を得ることができるのではないか？
\strut}\end{tabular}} \\ \hline
\multicolumn{1}{|l|}{A} & \multicolumn{1}{l|}{\begin{tabular}[c]{@{}l@{}}\parbox{37zw}{\strut{}
人が書いた文章を入力とした場合，高確率で人だと答えることができるが，締め切り直前にその実験結果を卒業研究論文に含めることができなかった．修正した現在のバージョンではその実験も追加した．
\strut}\end{tabular}} \\ \hline

\end{tabular}
\end{table}

\begin{table}[h]
\centering
\large
\begin{tabular}{ll}

& 大原　剛三　情報テクノロジー学科　教授 \\ \hline
\multicolumn{1}{|l|}{Q} & \multicolumn{1}{l|}{\begin{tabular}[c]{@{}l@{}}\parbox{37zw}{\strut{}
評価に使っている文章はほとんどが時事ニュースに関するものに見えるが，他の種類の文章は対象にしなくて良いのか？判定できている理由として，ChatGPT では学習が必要なために，最新の時事ニュースとの間にタイムラグがあるからだ，ということは考えられないか？
\strut}\end{tabular}} \\ \hline
\multicolumn{1}{|l|}{A} & \multicolumn{1}{l|}{\begin{tabular}[c]{@{}l@{}}\parbox{37zw}{\strut{}
タイムラグを利用しているわけではないと思うが，確かに他のクラスの文章での評価が必要．今後の課題とする．
\strut}\end{tabular}} \\ \hline

\end{tabular}
\end{table}

\begin{table}[h]
\centering
\large
\begin{tabular}{ll}

& 伊藤　弘大　情報テクノロジー学科　助教 \\ \hline
\multicolumn{1}{|l|}{Q} & \multicolumn{1}{l|}{\begin{tabular}[c]{@{}l@{}}\parbox{37zw}{\strut{}
Omission を実行する際にパラメータ $O$ があるが，いくつに設定しているか？
\strut}\end{tabular}} \\ \hline
\multicolumn{1}{|l|}{A} & \multicolumn{1}{l|}{\begin{tabular}[c]{@{}l@{}}\parbox{37zw}{\strut{}
3.14に設定されている．
\strut}\end{tabular}} \\ \hline

\end{tabular}
\end{table}

